\documentclass[10pt,a4paper]{amsart}
\usepackage[margin=19mm]{geometry}
\usepackage{amsmath,amssymb,mathtools}
\usepackage[scr=rsfs,cal=euler]{mathalfa}
\usepackage{xfrac}
\usepackage[unicode,hidelinks,bookmarks=false]{hyperref}
\newcommand{\ie}{i.e.\ }
\newcommand{\cf}{cf.\ }
\newcommand{\resp}{respectively\ }
\newcommand{\Subsection}{\subsection}
\providecommand{\nobreakdash}{\nobreak\hbox{-}}
\setlength{\emergencystretch}{2em}
\multlinegap0pt
\usepackage{tikz}
\usetikzlibrary{cd}
\usepackage{amssymb}
\newtheorem{prop}{Proposition}
\newcommand{\bA}{\mathbb{A}}
\newcommand{\bB}{\mathbb{B}}
\title{Homological algebra over non-unital rings and algebras, with applications to $(\infty, 1)$-categories}
\author{Eric Goubault, Eliot M\'edioni}
\date{Source: arXiv:2603.11937v1 (CC BY 4.0)}
\begin{document}
\maketitle

\noindent\emph{Source and licence.} Homological algebra over non-unital rings and algebras, with applications to $(\infty, 1)$-categories. Version arXiv:2603.11937v1 (CC BY 4.0), distributed by arXiv under the Creative Commons Attribution 4.0 International licence, http://creativecommons.org/licenses/by/4.0/, as recorded in the arXiv licence metadata for this exact version and shown on the abstract page https://arxiv.org/abs/2603.11937v1. Full licence notice: \texttt{licenses/arxiv-2603.11937-CC-BY-4.0.txt}.\par\smallskip

\noindent\textit{Editorial scope.} Counted unit: the article's prop prop (source0.tex lines 736--758). The article's complete original proof of it is delivered with it (source0.tex lines 760--838, one \texttt{proof} environment of 79 lines). No mathematics is written, repaired or re-derived: the statement and the proof are the article's own lines, in the article's own notation, and no retained line is substituted or re-flowed.  No further source-local context is needed: every cross-reference of every retained line resolves inside the two delivered spans.  Nothing was omitted from any retained span, so the package carries no omission record.  No retained line contains a citation, so the wrapper carries no bibliography and no bibliography entry.  The retained lines contain the article's own diagram code, which the wrapper typesets with the standard tikz packages; no image file is delivered and the package declares no asset dependency.  Editorial formatting only, in addition: an AMS \texttt{amsart} preamble that restates the article's own declarations and macros used by the retained lines, namely the environments \texttt{prop} on separate counters, together with the standard packages those lines need.  The retained environments print in this document as Proposition 1 in order of appearance, whereas the article prints the same items with the numbering of its own full document.  The complete native source of the article as distributed in the arXiv e-print for this exact version is bundled unchanged as \texttt{sources/196245/source0.tex}.
\begin{prop}
Let $g : \bA \to \bB$ be an $R$-algebra morphism.

Equip $\bB$ with the right $\bA$-module structure defined by
\[
\begin{cases}
b.a = b \times g(a) \text{ for } a\in \bA, b \in \bB 
\\
b.r = r.b\text{ for } r \in R, b \in \bB 
\end{cases}
\]

\begin{itemize}
    \item For any $\bA$-module $M$, $\bB \otimes_\bA M$ has a natural (left) $\bB$-module structure given by
    \[
    \begin{cases}
    r.(b \otimes m) = (r.b) \otimes m \\
    b' .(b \otimes m) = (b'b) \otimes m
    \end{cases}
    \]
    \item If $\bB$ is a firm $\bB$-module, then for any $\bA$-module $M$ we have that $\bB \otimes_\bA M$ is a firm $\bB$-module.
\end{itemize}
\end{prop}

\begin{proof}
%\todo{EG: \`a relire}
$\bullet$ The right $\bA$-module structure on $\bB$ is well-defined because $g$ is an $R$-algebra morphism. 
%\todo{5}
\iffalse
\begin{enumerate}
    \item $(b.a_1).a_2 = b g(a_1) g(a_2) 
(b.a_1).a_2= bg(a_1 a_1)$ so
\[(b.a_1).a_2= b.(a_1a_2)\]
    \item $(b.a).r = r. (b \times g(a)) = b \times ( r g(a)) $ by bilinearity and $r g(a) = g(ra) = g(a.r)$ so
    \[
    (b.a).r = b.(a.r)
    \]
    \item $(b.r).a = r.b \times g(a) = b \times g(a).r = (b.a).r$
    \item $b.(a_1 + a_2) = b g(a_1 + a_2) = bg(a_1) + bg(a_2)$ so
    \[b.(a_1 + a_2) = b.a_1 + b.a_2\]
\item The rest is trivial    
\end{enumerate}
\fi

$\bullet$ It is clear that $\bB \otimes_\bA M$ is indeed a left $\bB$-module.

$\bullet$ Taking either 
$$
U = \bB \otimes_\bB ( \bB \otimes_\bA M)
\text{ and }
p_U : (b,b',m) \mapsto b \otimes (b' \otimes m)
$$ 
or
$$
U = (\bB \otimes_\bB \bB ) \otimes_\bA M
\text{ and }
p_U : (b,b',m) \mapsto (b \otimes b') \otimes m
$$ such that $p$ is the canonical $R$-module homomorphism $bB \times  \bB \times M \to N \to U$, 
$(U, p_U)$ verifies the following universal property:

%\begin{center}
    \textit{For any $R$-module $N$, every $R$-trilinear map $f : bB \times  \bB \times M \to N$ such that
    \begin{itemize}
        \item $f(b.b' , b'', m) = f(b, b'.b'', m)$ for any $b, b',b'' \in \bB, m \in M$
        \item $f(b , b'.a, m) = f(b, b', a.m)$ for any $a \in \bA, b,b' \in \bB, m \in M$
    \end{itemize}
    there is a unique $R$-module homomorphism $\Tilde{f} : U \to N$ such that 
    \[ 
    \Tilde{f} = f \circ p_U
    \]}
%\end{center}

As such, there is a a canonical isomorphism 
\[
\phi : (\bB \otimes_\bB \bB) \otimes_\bA M 
\to \bB \otimes_\bB (\bB \otimes_\bA M) 
\]
such that 
$\phi ((b \otimes b') \otimes a )
= b \otimes (b' \otimes a)$
and then the universal maps
\[
\begin{cases}
\mu : \mu : \bB \otimes_\bB (\bB \otimes_\bA M)  \to (\bB \otimes_\bA M),
\quad
b \otimes (b' \otimes m) \mapsto (bb') \otimes m
\\
\nu :
\bB \otimes_\bB \bB  \to \bB,
\quad
b \otimes b' \mapsto (bb')
\end{cases}
\]
make the following diagram commute:
\[\begin{tikzcd}
	{(\bB \otimes_\bB \bB) \otimes_\bA M} & {\bB \otimes_\bB (\bB \otimes_\bA M)} \\
	& {\bB \otimes_\bA M}
	\arrow["\phi", from=1-1, to=1-2]
	\arrow["{\nu \otimes \mathrm{Id}_M}"', from=1-1, to=2-2]
	\arrow["\mu", from=1-2, to=2-2]
\end{tikzcd}\]
so if $\nu$ is an isomorphism, then $\mu$ is one too.\\
\end{proof}

\end{document}
